\chapter{Roles and permissions}
\label{ch:roles}

mybibli has three user roles\index{roles} ordered by privilege:

\begin{enumerate}
  \item \textbf{Anonymous}\index{anonymous} — anyone who lands on
        the site without a session cookie. Read-only access to the
        catalog. Cannot loan, edit or delete anything.
  \item \textbf{Librarian}\index{librarian} — authenticated user
        whose role is \texttt{librarian}. Day-to-day cataloging,
        loans, returns, location moves. Cannot manage users or
        change system settings.
  \item \textbf{Admin}\index{admin role} — authenticated user whose
        role is \texttt{admin}. Everything a Librarian can do, plus
        the \texttt{/admin} panel: users, reference data, trash,
        system settings.
\end{enumerate}

\section{Who can do what}

\begin{longtable}{@{} l c c c @{}}
\toprule
\textbf{Action} & \textbf{Anon} & \textbf{Lib} & \textbf{Admin} \\
\midrule
\endhead
Browse / search the catalog        & yes & yes & yes \\
View title and volume pages        & yes & yes & yes \\
View location tree                 & yes & yes & yes \\
\midrule
Scan / catalog a new title         & no  & yes & yes \\
Add or remove a volume             & no  & yes & yes \\
Move a volume between locations    & no  & yes & yes \\
Edit a title's metadata            & no  & yes & yes \\
Create / edit borrowers            & no  & yes & yes \\
Register / return loans            & no  & yes & yes \\
Edit storage location tree         & no  & yes & yes \\
\midrule
Open \texttt{/admin}               & no  & no  & yes \\
Create / deactivate users          & no  & no  & yes \\
Edit reference data                & no  & no  & yes \\
Restore from trash                 & no  & no  & yes \\
Permanently delete from trash      & no  & no  & yes \\
Edit system settings               & no  & no  & yes \\
\bottomrule
\end{longtable}

\section{Sessions and the inactivity timeout}

After login, mybibli sets a session cookie\index{session cookie}
that travels with every subsequent request. The cookie is:

\begin{itemize}
  \item \texttt{HttpOnly} — not exposed to JavaScript;
  \item \texttt{SameSite=Lax} — sent on top-level same-site
        navigations but not on cross-site requests;
  \item \texttt{Secure} (only when \texttt{MYBIBLI\_COOKIE\_SECURE}
        is set to \texttt{true} — see chapter~\ref{ch:installation}).
\end{itemize}

The session has an \textbf{inactivity timeout}\index{inactivity
timeout} (default 30 minutes, configurable from
\texttt{/admin > System > Loans}). Every authenticated request resets
the countdown; if no request lands within the timeout window, the
session is wiped server-side and the next request redirects to the
login page.

A small \emph{keep-alive toast} appears in the corner of the screen
roughly two minutes before expiry — click it to extend the session
without leaving your current page.

\section{The last-active-admin guard}

mybibli refuses to leave the installation without any active admin.
The guard applies to two actions:

\begin{itemize}
  \item \textbf{Deactivating a user.} If the user is the only
        active admin, the request returns a 409 conflict and the
        deactivation is refused.
  \item \textbf{Permanently deleting a user from the trash.} Same
        rule.
\end{itemize}

Self-deactivation is also blocked unconditionally (you cannot
deactivate the account you are currently logged in as).

\begin{warning}
The guard counts only \emph{active} admins. If you mass-deactivate
admins and the guard fires, undo by reactivating someone from the
trash tab. There is no recovery path through the wizard or the
database without manual SQL surgery.
\end{warning}

\section{Password reset}

mybibli does not implement password-reset emails — household /
small-community deployments rarely benefit from the email
plumbing. An admin can reset another user's password by:

\begin{enumerate}
  \item Open \texttt{/admin > Users}.
  \item Pick the user, click \textbf{Deactivate}.
  \item Click \textbf{Reactivate} from the trash tab.
  \item Re-create the user with a fresh password.
\end{enumerate}

A user who knows their own password can change it from the
\texttt{/account} page.

\section{Anonymous sessions}

Anonymous visitors also get a session row (used for CSRF protection
and pending HTMX updates). Anonymous sessions expire after seven
days of inactivity and are garbage-collected by a daily background
task — they do not affect logged-in users.
